\documentclass[11pt,a4paper]{article}
\usepackage{operationsform}
\begin{document}
\OPStitle{校正確認チェックリスト}{TPL-CHK-030}{0.2.0}
\OPSsection{1}{確認項目}
\OPScheck{1}{校正対象の原稿版・ゲラ版・媒体が特定されている}
\OPScheck{2}{誤字、脱字、記号、数値表記と名称の一致を検査した}
\OPScheck{3}{見出し、脚注、図表番号、リンク、字幕を確認した}
\OPScheck{4}{修正後の再校で反映を確認した}
\OPScheck{5}{未解決の表示不具合と責任者を記録した}
\OPScheck{6}{校正と科学的校閲の責務を区別した}
\OPSsection{2}{点検結果の記録}
\OPSblankshort{対象案件・対象版}
\OPSblankshort{担当者と点検日}
\OPSblanklong{未解決事項と対応}
\par\Needspace{4\baselineskip}\noindent\textbf{確認に用いた資料ID}\par\vspace{1mm}
{\fontsize{8.5}{13}\selectfont
\begin{tabularx}{\linewidth}{@{}XXXX@{}}
\textbf{対象・項目} & \textbf{内容・参照} & \textbf{担当・状態} & \textbf{備考}\\[2mm]
\strut & \strut & \strut & \strut\\[5mm]
\strut & \strut & \strut & \strut\\[5mm]
\end{tabularx}
}\vspace{2mm}\par
\end{document}
